% Sight Reduction Form + instructions
% Companion to the HO229 sight reduction program for the TI-85.
% Build:  pdflatex sight-reduction-form.tex
% Only stock LaTeX packages are used; no external fonts or files required.

\documentclass[10pt]{article}

\usepackage[letterpaper,top=0.6in,bottom=0.6in,left=0.7in,right=0.7in]{geometry}
\usepackage{array}
\usepackage{tabularx}
\usepackage{textcomp}
\usepackage{amsmath}
\usepackage{fancyhdr}
\usepackage{enumitem}
\usepackage{lastpage}

\pagestyle{fancy}
\fancyhf{}
\renewcommand{\headrulewidth}{0pt}
\fancyfoot[L]{\footnotesize Sight Reduction Form}
\fancyfoot[C]{\footnotesize\thepage\ of \pageref{LastPage}}
\fancyfoot[R]{\footnotesize Companion to \texttt{HO229} (TI-85)}

% ---------------------------------------------------------------- helpers
\newcommand{\dg}{\ensuremath{^\circ}}
\newcommand{\mn}{\ensuremath{'}}

% row strut: sets the writing height of every blank entry row
\newcommand{\rowstrut}{\rule[-1.0ex]{0pt}{3.6ex}}

% section heading bar
\newcommand{\secbar}[1]{%
  \noindent\fbox{\begin{minipage}{\dimexpr\textwidth-2\fboxsep-2\fboxrule\relax}%
  \textbf{#1}%
  \end{minipage}}\vspace{0.4ex}%
}

% one entry row: \fldrow{item no.}{description}
\newcommand{\fldrow}[2]{#1 & #2 & \rowstrut \\ \hline}
% a sub-item (indented, no item number)
\newcommand{\subrow}[1]{ & \hspace{1.2em}#1 & \rowstrut \\ \hline}
% a total row (bold)
\newcommand{\totrow}[1]{ & \hspace{1.2em}\textbf{#1} & \rowstrut \\ \hline}

% framed full-width note box (avoids nesting tabular inside a tabularx X cell)
\newcommand{\notebox}[1]{%
  \noindent\fbox{\begin{minipage}{\dimexpr\textwidth-2\fboxsep-2\fboxrule\relax}%
  #1%
  \end{minipage}}%
}

\newenvironment{formtable}{%
  \renewcommand{\arraystretch}{1.0}%
  \noindent\begin{tabular}{|c|p{4.7in}|p{1.6in}|}\hline
}{%
  \end{tabular}\vspace{0.8ex}%
}

\setlength{\parindent}{0pt}

% ================================================================
\begin{document}

% ------------------------------------------------ PAGE 1: THE FORM
\begin{center}
  {\Large\bfseries Sight Reduction Form}\\[0.3ex]
  {\footnotesize Vessel / Observer \rule[-0.5ex]{2.0in}{0.4pt}\qquad
   Sight No. \rule[-0.5ex]{0.7in}{0.4pt}}
\end{center}
\vspace{-0.6ex}

\secbar{Section I: Observations and Corrections}
\begin{formtable}
  \fldrow{1.}{Celestial Body}
  \fldrow{2.}{\textbf{Apparent Altitude} \quad Sextant Altitude (hs)}
  \subrow{$\pm$ Index Correction}
  \subrow{$-$ Dip Correction}
  \totrow{Total Apparent Altitude (ha)}
  \fldrow{3.}{\textbf{Observed Altitude} \quad $\pm$ Altitude Correction}
  \subrow{$\pm$ Additional Corrections (Atmospheric)}
  \subrow{$\pm$ Additional Corrections (Mars / Venus / Moon)}
  \totrow{Total Observed Altitude (ho)}
\end{formtable}

\secbar{Section II: Time and Dead Reckoning}
\begin{formtable}
  \fldrow{4.}{Date (GMT)}
  \fldrow{5.}{DR Latitude}
  \fldrow{6.}{DR Longitude}
  \fldrow{7.}{Time (GMT)}
\end{formtable}

\secbar{Section III: Local Hour Angle and Declination}
\begin{formtable}
  \fldrow{8.}{\textbf{GHA} \quad Tabulated GHA}
  \subrow{$+$ GHA Increment}
  \subrow{$+$ SHA (stars) \emph{or} v-correction (Moon / planets, when applicable)}
  \totrow{Total GHA}
  \fldrow{9.}{\textbf{LHA} \quad (a) Choose Assumed Longitude so LHA minutes $=00.0\mn$}
  \subrow{{\footnotesize W: a-minutes $=$ GHA minutes, LHA $=$ GHA$-$a. \quad E: a-minutes $=60.0\mn-$GHA minutes, LHA $=$ GHA$+$a.}}
  \subrow{(b) Normalize: $+360$\dg{} if LHA $<0$\dg; $-360$\dg{} if LHA $\geq360$\dg}
  \totrow{Total LHA \quad {\footnotesize (record normalized value, $0$\dg{} to $<360$\dg)}}
  \fldrow{10.}{\textbf{Declination} \quad Tabulated Dec. \quad {\footnotesize (d-corr'n factor: \rule[-0.4ex]{0.7in}{0.4pt})}}
  \subrow{$\pm$ d-correction}
  \totrow{Total Declination (Dec.)}
  \fldrow{11.}{\textbf{Assumed Latitude} \quad Same / Contrary (compared to Dec.\ hemisphere)}
\end{formtable}

\secbar{Section IV: Determining a Line of Position}
\begin{formtable}
  \fldrow{12.}{\textbf{Computed Altitude} \quad Tabulated hc \quad {\footnotesize (d: \rule[-0.4ex]{0.7in}{0.4pt})}}
  \subrow{$\pm$ Declination Increment (Dec.\ minutes $/\,60\times$ d)}
  \totrow{Total Computed Altitude (hc)}
  \fldrow{13.}{\textbf{Altitude Intercept} \quad (ho or hc) whichever is \emph{larger} \quad {\footnotesize ho $=$ I.3}}
  \subrow{$-$ (ho or hc) whichever is \emph{smaller} \quad {\footnotesize hc $=$ IV.12}}
  \totrow{( T / A ) Intercept \quad {\footnotesize (Toward if ho $>$ hc, Away if ho $<$ hc)}}
  \fldrow{14.}{\textbf{Azimuth Angle (Z)}}
  \fldrow{15.}{\textbf{Azimuth (Zn)}}
\end{formtable}

\vspace{-0.2ex}
\notebox{%
\textbf{Zn from Z --- quadrant rules}\\[0.6ex]
\begin{tabular}{@{}l@{\hspace{1.5em}}l@{\hspace{2.5em}}l@{}}
\textbf{N Lat.:} & LHA $>180$\dg: \ Zn $=$ Z & LHA $<180$\dg: \ Zn $=360$\dg$-$Z\\[0.2ex]
\textbf{S Lat.:} & LHA $>180$\dg: \ Zn $=180$\dg$-$Z & LHA $<180$\dg: \ Zn $=180$\dg$+$Z\\
\end{tabular}}


\clearpage
\newgeometry{top=0.45in,bottom=0.45in,left=0.65in,right=0.65in}
\small\linespread{0.95}\selectfont

% ------------------------------------- PAGE 2: HOW TO FILL IT OUT
\begin{center}
  {\large\bfseries How to Fill Out the Sight Reduction Form}
\end{center}
\vspace{0.4ex}

The form reduces a sextant reading to a \emph{line of position} (LOP): a line
somewhere on which the observer must lie. It compares the altitude you actually
measured (\textbf{ho}) against the altitude you \emph{would have} measured from an
assumed position (\textbf{hc}). That difference, in nautical miles, is the
\textbf{intercept}; the direction to the body is the \textbf{azimuth}. One sight
gives one LOP; two or more crossing LOPs give a fix.

\vspace{0.5ex}
\textbf{Angle notation.} Enter angles as degrees and decimal minutes
(e.g.\ $45\dg\,13.4\mn$). \texttt{HO229} takes the same thing compressed to
\texttt{DD.MMm}: $45\dg\,13.4\mn$ is keyed \texttt{45.134}; LHA as \texttt{DDD.MMm}.

\vspace{0.8ex}
\textbf{Section I --- Observations and Corrections}
\begin{description}[leftmargin=2.4em,style=nextline,itemsep=0.1ex,topsep=0.2ex,parsep=0pt]
\item[1. Celestial body.] Sun (specify lower or upper limb), Moon, a planet, or a
  named star. This determines which correction tables apply later.
\item[2. Apparent altitude.] Start with the raw sextant altitude \textbf{hs}.
  Apply the \emph{index correction} for your sextant's own error (on the arc is
  subtractive, off the arc additive), then \emph{dip} for your height of eye,
  which is always subtractive. The result is \textbf{ha}.
\item[3. Observed altitude.] Apply the main altitude correction --- refraction,
  semi-diameter, parallax --- from the almanac's critical tables, entered with ha.
  Add the atmospheric correction if temperature or pressure are non-standard, and
  the extra correction for Mars, Venus, or the Moon where it applies. The result,
  \textbf{ho}, is the true altitude of the body's centre above the celestial
  horizon.
\end{description}

\textbf{Section II --- Time and Dead Reckoning}
\begin{description}[leftmargin=2.4em,style=nextline,itemsep=0.1ex,topsep=0.2ex,parsep=0pt]
\item[4--7. Date, DR position, time.] Record GMT (UT) to the second --- four seconds
  of time is a nautical mile of longitude at the equator. The DR position is your
  best estimate at the moment of the sight; you round it to get the assumed
  position in items 9 and 11.
\end{description}

\textbf{Section III --- Local Hour Angle and Declination}
\begin{description}[leftmargin=2.4em,style=nextline,itemsep=0.1ex,topsep=0.2ex,parsep=0pt]
\item[8. GHA.] Take the tabulated GHA for the whole hour from the daily pages, add
  the increment for the minutes and seconds. For a star, add its SHA to the GHA of
  Aries; for the Moon, apply the v-correction.
\item[9. LHA.] Choose the assumed longitude so the resulting LHA has $00.0\mn$ of
  arc. West: give it the same minutes as GHA, then LHA $=$ GHA $-$ a. East: give it
  $60.0\mn-$GHA minutes, then LHA $=$ GHA $+$ a. Normalize afterwards --- add
  $360\dg$ if negative, subtract $360\dg$ if $360\dg$ or more. W and E are
  conditions telling you which operation to perform, not quantities to solve for.
\item[10. Declination.] Tabulated declination for the hour, plus the d-correction
  for the minutes. Note carefully whether it is North or South.
\item[11. Assumed latitude.] Round the DR latitude to a whole degree. Mark
  \emph{Same} if the assumed latitude and the declination are in the same
  hemisphere, \emph{Contrary} if opposite --- the tables are laid out in those two
  cases and it is an easy place to go wrong.
\end{description}

\textbf{Section IV --- Determining a Line of Position}
\begin{description}[leftmargin=2.4em,style=nextline,itemsep=0.1ex,topsep=0.2ex,parsep=0pt]
\item[12. Computed altitude.] Enter Pub.\ 229 with LHA, assumed latitude, and
  declination; take out tabulated hc, d, and Z. Apply the declination increment
  (the odd minutes of declination, times d, divided by 60) to get \textbf{hc}.
\item[13. Altitude intercept.] Subtract the smaller of ho and hc from the larger;
  one minute of arc is one nautical mile. Label it \textbf{T} (Toward) if ho is
  greater, \textbf{A} (Away) if hc is greater --- \emph{Ho Mo To}, ho more, toward.
\item[14. Azimuth angle Z.] From the same table entry. Z is measured from the
  \emph{elevated pole} --- the pole nearer to you --- and runs $0\dg$ to $180\dg$
  east or west.
\item[15. Azimuth Zn.] Convert Z to a true bearing $0\dg$--$360\dg$ using the
  quadrant rules boxed at the foot of page 1. This is the step the rules exist for:
  Z alone does not say which side of the meridian the body lies on, and in southern
  latitudes it is measured from the south pole rather than the north.
\end{description}

\textbf{Plotting.} From the assumed position, draw the azimuth line along Zn. Step
off the intercept along it --- toward the body or away as labelled --- and mark that
point. The LOP is the perpendicular to the azimuth through that point. Repeat with a
second body; the crossing is your fix.

\vspace{0.8ex}
\hrule
\vspace{0.6ex}

\textbf{Using the \texttt{HO229} program in place of the tables}

\vspace{0.4ex}
The program computes items 12 and 14 directly and does item 15 for you. Run it and
enter, in \texttt{DD.MMm} notation:

\begin{description}[leftmargin=2.4em,style=nextline,itemsep=0pt,topsep=0.2ex,parsep=0pt]
\item[\texttt{LAT N+ S-}] the assumed latitude from item 11; negative for South.
\item[\texttt{LHA DDD.MMm}] the total LHA from item 9.
\item[\texttt{DEC N+ S-}] the total declination from item 10; negative for South.
\end{description}

It returns \texttt{HC}, \texttt{Z}, and \texttt{ZN} in decimal degrees.
\texttt{HC} goes straight into item 12 --- there is no declination increment to
apply, because the program uses the exact declination rather than a tabulated
whole degree. \texttt{Z} is reported from the elevated pole, in the same sense as
Pub.\ 229, so it goes in item 14 and the quadrant rules on page 1 apply to it
unchanged. \texttt{ZN} is item 15, already converted --- work item 15 by hand
anyway the first few times and check that the two agree.

\vspace{0.6ex}
Because the program needs no whole-degree entering arguments, you may use your actual
DR position and exact LHA instead of an assumed position --- but then plot from the
DR, consistently across every sight in a round. Sections I--III are unchanged: the
program replaces the table lookup, and nothing else.

\end{document}
